\documentclass[11pt,a4paper]{article}

% Appunti personali di lettura: non sono un testo da presentare.
\usepackage[utf8]{inputenc}
\usepackage[T1]{fontenc}
\usepackage[italian]{babel}
\usepackage{amsmath,amssymb,mathtools}
\usepackage[a4paper,margin=2.45cm]{geometry}
\usepackage{enumitem}
\usepackage{booktabs}
\usepackage{xcolor}
\usepackage{tikz}
\usetikzlibrary{arrows.meta,positioning,calc}
\usepackage{hyperref}

\title{Appunti di lettura\\\large Durante, Fernández-Sánchez e Pappadà (2015), ``Copulas, diagonals, and tail dependence''}
\author{Appunti personali di Massimo}
\date{}

\newcommand{\R}{\mathbb{R}}
\newcommand{\Prob}{\mathbb{P}}
\newcommand{\one}{\boldsymbol{1}}
\newcommand{\Idea}{\paragraph{Idea.}}
\newcommand{\Esempio}{\paragraph{Esempio.}}
\newcommand{\Ricorda}{\paragraph{Da ricordare.}}
\newcommand{\Salta}{\paragraph{Cosa puoi saltare per ora.}}

\begin{document}
\maketitle

\paragraph{Come usare questi appunti.}
Il paper di Durante et al. è il terzo passo dopo Schmidt e Hua--Joe. Questi
appunti sono scritti per essere ripetuti a voce, non per essere consegnati: ogni
concetto parte da una domanda concreta, poi riceve il suo nome tecnico e infine
la formula. Le figure sono diagrammi didattici, non dati reali. Userò ``TDC''
per il coefficiente asintotico e ``TCF'' per la funzione a soglia finita: la
differenza tra i due è la cosa più importante del paper.

\section{La domanda da cui partire}

Immagina di osservare ogni giorno le perdite di due banche. Una domanda
finanziaria naturale è: ``quanto spesso perdono entrambe nello stesso giorno?''
Ma ``spesso'' può voler dire due cose diverse. Possiamo chiederlo solo per
perdite infinitamente estreme, oppure per perdite molto grandi ma ancora
osservabili nel campione.

Schmidt risponde alla prima domanda con il tail-dependence coefficient (TDC),
un limite. Hua--Joe aggiunge una misura della velocità con cui quel limite viene
raggiunto. Durante, Fernández-Sánchez e Pappadà propongono di guardare tutta la
curva della dipendenza nella regione finita della coda. La chiamano \emph{tail
concentration function}, o TCF.

\Idea Il TDC è come conoscere solo l'ultima fermata di un viaggio. La TCF
mostra il percorso: due modelli possono arrivare allo stesso punto finale, ma
uno può restare molto più rischioso per tutte le soglie che riusciamo davvero a
osservare.

\section{Percentili, copula e sezione diagonale}

Prendiamo una coppia di rendimenti $(X,Y)$. Una singola osservazione è una
coppia di numeri, per esempio $(0.02,-0.01)$. Prima di confrontare i due titoli
trasformiamo ciascun numero in un percentile: $U=F_X(X)$ e $V=F_Y(Y)$. Un
percentile vicino a $0$ significa ``molto basso rispetto alla propria serie'';
uno vicino a $1$ significa ``molto alto''.

La copula è la distribuzione congiunta di $(U,V)$. In parole semplici, le
marginali dicono quanto è rischioso ciascun titolo preso da solo; la copula
dice come i loro percentili si combinano. Per marginali continue:
\begin{equation}
 F(x,y)=C\bigl(F_X(x),F_Y(y)\bigr).
 \label{eq:sklar}
\end{equation}

\begin{figure}[ht]
\centering
\begin{tikzpicture}[x=4.0cm,y=4.0cm,font=\small]
  \draw[->] (0,0)--(1.08,0) node[right] {$u$ (percentile di $X$)};
  \draw[->] (0,0)--(0,1.08) node[above] {$v$ (percentile di $Y$)};
  \draw[gray] (0,0) rectangle (1,1);
  \draw[dashed,gray] (0.25,0)--(0.25,1);
  \draw[dashed,gray] (0,0.25)--(1,0.25);
  \fill[orange!20] (0,0) rectangle (0.25,0.25);
  \node[align=center] at (0.13,0.13) {entrambi\\molto bassi};
  \foreach \x/\y in {0.12/0.18,0.19/0.10,0.34/0.40,0.45/0.38,0.57/0.62,0.66/0.71,0.77/0.74,0.91/0.88}{\fill[blue!70] (\x,\y) circle (1.8pt);}
  \fill[red!75] (0.19,0.10) circle (2.8pt);
  \node[align=left,anchor=west] at (0.42,1.03) {ogni punto è una coppia\\di percentili osservata};
\end{tikzpicture}
\caption{Il quadrato dei percentili: la zona arancione è la lower tail. La diagonale $u=v$ confronta soglie uguali nelle due marginali. Schema concettuale.}
\label{fig:percentili}
\end{figure}

La \emph{sezione diagonale} è la curva ottenuta guardando solo la linea
$u=v$:
\begin{equation}
 \delta_C(t)=C(t,t)=\Prob(U\leq t,V\leq t).
 \label{eq:diagonal}
\end{equation}
\Esempio Se $t=0.05$, $\delta_C(0.05)$ è la probabilità che entrambi i
percentili siano sotto il quinto percentile. È un conteggio nella zona
lower-left del quadrato. Per la upper tail useremo la probabilità che entrambi
superino una soglia alta; questa viene descritta dalla survival copula
$\overline C$.

\Salta per ora i teoremi che caratterizzano tutte le funzioni possibili come
diagonali di una copula. Il paper mostra anche due costruzioni, diagonal copula
e Bertino copula, per ricordare una cosa semplice: la stessa diagonale può
appartenere a molte copule diverse. La diagonale è informativa sulle code, ma
non ricostruisce ogni dettaglio della dipendenza.

\section{Il limite TDC non racconta tutto}

Se ci interessa la lower tail, Schmidt definisce
\begin{equation}
 \lambda_L=\lim_{t\downarrow0}\frac{\delta_C(t)}{t}.
 \label{eq:lambda-lower}
\end{equation}
Il numeratore è la probabilità che entrambi siano sotto $t$; il denominatore
$t$ è la probabilità che un singolo percentile sia sotto $t$. Il rapporto è
quindi una quota condizionata, scritta senza dover scegliere quale variabile
condizionare.

Per la upper tail:
\begin{equation}
 \lambda_U=\lim_{t\uparrow1}\frac{1-2t+\delta_C(t)}{1-t}.
 \label{eq:lambda-upper}
\end{equation}
Qui $1-2t+\delta_C(t)$ è la probabilità che entrambi superino $t$, mentre
$1-t$ è la probabilità di un singolo superamento.

\Ricorda Il TDC è il valore al bordo: dice cosa accade quando $t$ va a $0$ o
a $1$. Non dice se, per esempio, a $t=0.05$ la probabilità congiunta sia
ancora molto alta. Per un campione finanziario la differenza tra limite e
soglia osservabile può essere decisiva.

\section{La TCF: costruirla e leggerla punto per punto}

Gli autori mettono insieme le due frazioni in una sola curva:
\begin{equation}
 q_C(t)=
 \begin{cases}
  \displaystyle\frac{\delta_C(t)}{t}, & 0<t\leq0.5,\\[1.1ex]
  \displaystyle\frac{1-2t+\delta_C(t)}{1-t}, & 0.5<t<1.
 \end{cases}
 \label{eq:tcf}
\end{equation}

\paragraph{Perché dividiamo?}
Supponiamo $t=0.05$ e che $2$ osservazioni su $100$ siano entrambe nella
lower tail. La probabilità congiunta empirica è $0.02$; dividendo per $0.05$
otteniamo $q_C(0.05)=0.4$. La domanda diventa: ``tra le occasioni in cui un
singolo titolo è sotto il quinto percentile, quanto è grande la quota in cui
anche l'altro lo è?''

Per $t>0.5$ il ragionamento è speculare: contiamo le osservazioni sopra $t$ e
dividiamo per $1-t$, la probabilità marginale di superare la soglia. A $t=0.5
$ la curva passa dalla metà inferiore alla metà superiore. Il valore centrale
è legato alla misura di associazione di Blomqvist:
\begin{equation}
 q_C(0.5)=\frac{1+\beta_C}{2}.
 \label{eq:blomqvist}
\end{equation}
In pratica, $\beta_C$ è un indice di quanta massa cade nei due angoli rispetto
alla parte centrale del quadrato.

\begin{figure}[ht]
\centering
\begin{tikzpicture}[x=1.05cm,y=0.75cm,font=\small]
  \draw[->] (0,0)--(10.2,0) node[right] {$t$};
  \draw[->] (0,0)--(0,4.3) node[above] {$q_C(t)$};
  \draw[dashed,gray] (5,0)--(5,4) node[above] {$0.5$};
  \draw[blue!70,thick] (0.35,0.35) .. controls (1.6,0.45) and (3.6,1.10) .. (5,1.8)
    .. controls (6.4,2.25) and (8.6,2.65) .. (9.8,2.95);
  \draw[red!70,thick] (0.35,0.10) .. controls (1.6,0.20) and (3.6,0.60) .. (5,1.8)
    .. controls (6.4,2.90) and (8.6,3.55) .. (9.8,3.90);
  \node[blue!70,align=center] at (2.0,2.0) {TCF più bassa\\nella lower tail};
  \node[red!70,align=center] at (7.5,1.40) {più rischio a\\soglia finita};
  \node[align=center] at (5,-0.75) {limite a sinistra $=\lambda_L$; limite a destra $=\lambda_U$};
\end{tikzpicture}
\caption{Una TCF si legge da sinistra a destra: lower tail, centro, upper tail. Curve alte o lente indicano più dipendenza a soglia finita. Disegno schematico.}
\label{fig:tcf-reading}
\end{figure}

Una curva che scende lentamente verso il suo limite resta alta per molte
soglie: in finanza significa più probabilità di perdite congiunte osservabili.
Il paper mostra, ad esempio, che Gaussian, Frank e Plackett hanno TDC nullo
nella parte confrontata, ma la Gaussiana convergente più lentamente. Quindi il
limite non basta a ordinare il rischio finito.

\section{Le famiglie: quale coda mette in evidenza ciascuna?}

È più utile leggere le famiglie come scelte di comportamento che come elenco di
nomi.

\paragraph{Gaussian.} Immagina una nuvola con relazione centrale regolare. La
correlazione può essere positiva, ma la probabilità di estremi congiunti tende
a zero: TDC nullo. La TCF può però rimanere alta a soglie finite e scendere
lenta. È un buon riferimento per ``dipendenza senza vera tail dependence''.

\paragraph{Frank e Plackett.} Sono famiglie che possono descrivere dipendenza
centrale, ma nella figura del paper si spengono più rapidamente nelle code.
Sono utili come confronto: stesso centro, meno persistenza estrema.

\paragraph{Clayton.} Mette più dipendenza nella lower tail. Se interpretiamo
valori bassi come perdite, la curva lower resta alta più a lungo: a soglia
finita il rischio congiunto è relativamente conservativo. Il suo upper TDC è
nullo.

\paragraph{Survival Clayton.} È la versione riflessa: il comportamento forte è
nella upper tail invece che nella lower tail. Serve a ricordare che una famiglia
può essere asimmetrica tra le due code.

\paragraph{Gumbel e Galambos.} Entrambe privilegiano la upper tail nel confronto
del paper; le TCF risultano molto vicine. Non significa che le copule siano
identiche in ogni punto, solo che la loro descrizione diagonale delle code è
simile nel grafico.

\paragraph{Student-$t$.} Le code marginali sono pesanti e la dipendenza di coda
è positiva in entrambe le direzioni, in modo simmetrico nel caso considerato.
Con $\nu=4$ gradi di libertà è un modello naturale quando gli estremi sono più
frequenti di quanto preveda la normale.

\Salta per ora le formule specifiche dei parametri. Il messaggio da ripetere è
che la famiglia decide dove mettere la dipendenza: lower, upper, entrambe o
nessuna in senso asintotico. La TCF aiuta a scegliere anche quando i TDC delle
candidate coincidono.

\section{Come stimare la TCF: contare punti, non magie}

Supponiamo di avere $n$ coppie $(X_i,Y_i)$. Trasformiamo ciascuna coordinata in
un rango o percentile empirico $(U_i,V_i)$. Per un valore $t$, guardiamo il
quadrato $[0,t]^2$ e contiamo quanti punti cadono dentro:
\begin{equation}
 C_n(t,t)=\frac1n\sum_{i=1}^n
 \mathbf{1}_{\{U_i\leq t,\,V_i\leq t\}}.
 \end{equation}
Poi dividiamo per $t$ nella lower tail. Nella upper tail contiamo i punti sopra
le due soglie e dividiamo per $1-t$. È esattamente la formula TCF, solo con
conteggi osservati al posto della copula teorica.

\Esempio Con $n=20$ e $t=0.2$, il quadrato lower contiene $4$ posizioni
marginali. Se due punti cadono davvero nell'angolo lower-left, la stima è
$q_{\mathrm{emp}}(0.2)=2/20\,/\,0.2=0.5$. Se non cade alcun punto, la stima
è zero, ma non sappiamo se la dipendenza sia davvero assente: potremmo avere
semplicemente un campione povero di estremi.

\begin{figure}[ht]
\centering
\begin{tikzpicture}[x=4.1cm,y=4.1cm,font=\small]
  \draw[->] (0,0)--(1.08,0) node[right] {$U$};
  \draw[->] (0,0)--(0,1.08) node[above] {$V$};
  \draw[gray] (0,0) rectangle (1,1);
  \draw[dashed,orange!80,thick] (0.2,0)--(0.2,1);
  \draw[dashed,orange!80,thick] (0,0.2)--(1,0.2);
  \fill[orange!25] (0,0) rectangle (0.2,0.2);
  \foreach \x/\y in {0.08/0.12,0.15/0.07,0.12/0.29,0.31/0.35,0.45/0.55,0.60/0.65,0.72/0.80,0.85/0.77,0.93/0.91}{\fill[blue!65] (\x,\y) circle (1.9pt);}
  \fill[red!75] (0.08,0.12) circle (3pt);
  \fill[red!75] (0.15,0.07) circle (3pt);
  \node[align=center] at (0.10,0.10) {};
  \node[align=left,anchor=west] at (0.29,0.16) {2 punti su 20\\$q_{\rm emp}(0.2)=0.5$};
  \node[align=center] at (0.5,-0.20) {la soglia seleziona un conteggio};
\end{tikzpicture}
\caption{Stima lower della TCF: si contano i punti nel quadrato selezionato e si normalizza per la probabilità marginale della soglia. Schema didattico.}
\label{fig:tcf-count}
\end{figure}

Il problema pratico è vicino ai bordi. Se $t$ è troppo piccolo, $[0,t]^2$
contiene pochissimi punti e la curva diventa rumorosa. Se $t$ è troppo grande,
non stiamo più misurando una regione veramente estrema. Per questo il paper
consiglia di tracciare la TCF empirica solo su $[\varepsilon,1-\varepsilon]$,
con $\varepsilon$ scelto in relazione a $n$.

In formula, usando le distribuzioni empiriche $F_n,G_n$ e la copula empirica
$C_n$:
\begin{equation}
 q_{\mathrm{emp}}(t)=
 \frac{C_n(t,t)}{t}\mathbf{1}_{(0,0.5]}(t)+
 \frac{1-2t+C_n(t,t)}{1-t}\mathbf{1}_{(0.5,1)}(t).
 \label{eq:tcf-emp}
\end{equation}
La TCF empirica dipende quindi dal campione: non va confusa con il TDC
teorico di una famiglia.

\section{Scegliere una copula: la mappa costruita dalle distanze}

Supponiamo di avere quattro modelli candidati. Vorremmo una risposta visiva a
``quale curva assomiglia di più a quella osservata?''

\paragraph{Prima analogia: una mappa.} Se conosco le distanze tra città ma non
le loro coordinate, posso costruire una mappa approssimata. La posizione di due
città vicine non è una nuova variabile: è un modo di visualizzare le distanze.
Il paper fa la stessa cosa con le TCF.

La distanza tra TCF empirica e modello $C_i$ è
\begin{equation}
 d(C_n,C_i)=\left[\int_0^1
 \bigl(q_{\mathrm{emp}}(t)-q_{C_i}(t)\bigr)^2\,dt\right]^{1/2}.
 \label{eq:tcf-distance}
\end{equation}
È una distanza integrata: misura quanto le due curve si separano in media lungo
le soglie. Si calcolano anche le distanze tra ogni coppia di modelli.

\paragraph{Solo dopo arriva il nome tecnico.} La matrice di tutte queste
distanze viene trasformata in un piano con la \emph{Multidimensional Scaling}
(MDS). Il piano è una mappa approssimata; non è una nuova misura statistica e
non sostituisce la curva originale. Se due punti sono vicini, i rispettivi
modelli hanno TCF simili secondo la distanza scelta.

\begin{figure}[ht]
\centering
\begin{tikzpicture}[font=\small,>=Latex]
  \begin{scope}[xshift=0cm]
    \draw[->] (0,0)--(3.0,0) node[right] {$t$};
    \draw[->] (0,0)--(0,2.5) node[above] {$q(t)$};
    \draw[blue!70,thick] (0.2,0.3)..controls(0.9,0.45)and(1.8,1.2)..(2.8,1.9);
    \draw[red!70,thick] (0.2,0.2)..controls(0.9,0.7)and(1.8,1.0)..(2.8,2.1);
    \node[align=center] at (1.5,-0.45) {TCF empirica\\e TCF di un modello};
  \end{scope}
  \draw[->,thick] (3.4,1.1)--(4.4,1.1) node[midway,above] {distanze};
  \begin{scope}[xshift=5cm]
    \draw[->] (0,0)--(3.0,0) node[right] {MDS 1};
    \draw[->] (0,0)--(0,2.5) node[above] {MDS 2};
    \fill[red!75] (1.4,1.25) circle (3pt) node[above right] {$C_n$};
    \fill[blue!70] (1.1,1.4) circle (2.3pt) node[above left] {Clayton};
    \fill[blue!70] (1.8,1.1) circle (2.3pt) node[below right] {Gumbel};
    \fill[blue!70] (2.6,2.0) circle (2.3pt) node[above] {Gaussian};
    \draw[dashed,gray] (1.4,1.25)--(1.1,1.4);
    \draw[dashed,gray] (1.4,1.25)--(1.8,1.1);
    \draw[dashed,gray] (1.4,1.25)--(2.6,2.0);
    \node[align=center] at (1.5,-0.45) {modelli vicini\\hanno TCF simili};
  \end{scope}
\end{tikzpicture}
\caption{Dalla differenza tra curve alla mappa MDS. La mappa aiuta a vedere, ma resta una rappresentazione approssimata.}
\label{fig:mds}
\end{figure}

Il paper usa simulazioni con $n=250$ e $500$, copule Clayton, Gumbel e Gaussian
e tau di Kendall $0.25,0.5,0.75$. Clayton viene riconosciuta abbastanza bene;
Gumbel è più difficile per associazioni basse; la Gaussiana migliora con
campione grande e tau almeno $0.5$. È una verifica preliminare del metodo, non
una garanzia universale.

\Salta per ora il dettaglio della funzione di stress. Devi solo sapere che
controlla quanto la mappa bidimensionale distorce le distanze; il paper indica
uno stress sotto circa $2.5\%$ come rappresentazione pratica.

\section{Clustering: quali mercati perdono insieme?}

Con molte serie finanziarie la domanda cambia: ``quali mercati hanno pattern
simili quando le perdite diventano grandi?'' Non vogliamo raggruppare soltanto
mercati con la stessa scala o volatilità. Vogliamo raggruppare quelli che
cadono insieme nella lower tail.

\subsection{Prima togliamo scala e volatilità}
Ogni serie ha dinamica e volatilità proprie. Per esempio, una serie può avere
periodi di calma e crisi, indipendentemente dalla copula con un'altra. Gli
autori procedono così:
\begin{enumerate}[leftmargin=*]
  \item adattano a ogni rendimento un modello ARMA--GARCH: ARMA descrive la
  dipendenza temporale media, GARCH la volatilità che cambia nel tempo;
  \item usano innovazioni Student-$t$ per permettere code marginali pesanti;
  \item trasformano i residui con la loro distribuzione stimata in valori
  uniformi, chiamati probability integral transforms (PIT);
  \item applicano la copula ai PIT. La copula misura il legame tra serie dopo
  aver tolto scala, media e volatilità marginale.
\end{enumerate}

\paragraph{Detto senza sigle.} Prima rendo confrontabili le serie, poi guardo
solo se i loro percentili estremi si presentano insieme. ARMA--GARCH e Student-$t$
sono strumenti per descrivere ciascuna marginale; la TCF è lo strumento per la
dipendenza tra marginali.

\subsection{Una distanza che guarda soltanto le perdite}
Per ogni coppia di serie $i,j$ si calcola la TCF empirica lower ai punti
$t=k/T$, con $k=1,\ldots,T/2$. La copula comonotona, che muove sempre due
variabili insieme, avrebbe TCF lower uguale a $1$. Gli autori definiscono:
\begin{equation}
 d_{ij}=\sum_{k=1}^{T/2}
 \left[q_{\mathrm{emp}}^{ij}\!\left(\frac{k}{T}\right)-1\right]^2.
 \label{eq:cluster-distance}
\end{equation}
Se $d_{ij}=0$, le due serie si comportano come perfettamente associate nella
parte osservata della lower tail. Un valore grande indica che non condividono
lo stesso pattern di perdite congiunte.

\begin{figure}[ht]
\centering
\begin{tikzpicture}[font=\scriptsize,>=Latex]
  \begin{scope}[xshift=0cm]
    \draw[rounded corners,fill=blue!5] (0,0) rectangle (2.4,2.3);
    \node[align=center] at (1.2,1.95) {rendimenti\\grezzi};
    \draw[blue!60] (0.2,0.45)--(0.65,0.75)--(1.0,0.55)--(1.45,1.2)--(2.2,0.85);
    \node[align=center] at (1.2,0.22) {scale diverse};
  \end{scope}
  \draw[->,thick] (2.8,1.15)--(3.8,1.15) node[midway,above] {modelli marginali};
  \begin{scope}[xshift=4.2cm]
    \draw[rounded corners,fill=green!7] (0,0) rectangle (2.4,2.3);
    \node[align=center] at (1.2,1.95) {ranghi/PIT};
    \draw[green!60!black] (0.2,0.5)--(0.65,0.9)--(1.0,0.55)--(1.45,1.35)--(2.2,0.95);
    \node[align=center] at (1.2,0.22) {percentili confrontabili};
  \end{scope}
  \draw[->,thick] (6.9,1.15)--(7.9,1.15) node[midway,above] {TCF lower};
  \begin{scope}[xshift=8.3cm]
    \draw[rounded corners,fill=orange!8] (0,0) rectangle (2.4,2.3);
    \node[align=center] at (1.2,1.95) {matrice\\$d_{ij}$};
    \draw[fill=orange!35] (0.35,0.35) rectangle (0.85,0.85);
    \draw[fill=orange!55] (0.95,0.35) rectangle (1.45,0.85);
    \draw[fill=orange!75] (1.55,0.35) rectangle (2.05,0.85);
    \draw[fill=orange!55] (0.35,0.95) rectangle (0.85,1.45);
    \draw[fill=orange!35] (0.95,0.95) rectangle (1.45,1.45);
    \draw[fill=orange!75] (1.55,0.95) rectangle (2.05,1.45);
    \node[align=center] at (1.2,0.15) {somiglianze lower-tail};
  \end{scope}
  \draw[->,thick] (11.0,1.15)--(12.0,1.15) node[midway,above] {clustering};
  \begin{scope}[xshift=12.3cm]
    \draw[rounded corners,fill=red!7] (0,0) rectangle (2.4,2.3);
    \node[align=center] at (1.2,1.95) {cluster};
    \node[fill=red!20,rounded corners] at (0.7,1.25) {Europa};
    \node[fill=red!20,rounded corners] at (1.7,1.25) {USA};
    \node[fill=red!20,rounded corners] at (1.2,0.55) {Pacifico};
  \end{scope}
\end{tikzpicture}
\caption{Percorso concettuale del clustering: dai rendimenti si eliminano gli effetti marginali, poi si confrontano le TCF lower e infine si formano i gruppi.}
\label{fig:pipeline-cluster}
\end{figure}

La matrice $(d_{ij})$ può essere passata a un clustering gerarchico
agglomerativo. Single, complete e average linkage sono tre modi diversi di
definire la distanza tra gruppi; il paper usa complete linkage nel suo esempio
MSCI.

\subsection{Esempio MSCI e diversificazione}
Il dataset contiene 23 indici MSCI di mercati sviluppati, rendimenti giornalieri
dal 4 giugno 2002 al 10 giugno 2010, con $T=2093$ osservazioni. Il dendrogramma
ottenuto con complete linkage suggerisce $4$ cluster. La lettura è soprattutto
geografica: mercati europei e scandinavi, un gruppo nordamericano e gruppi
dell'area del Pacifico.

Se due mercati stanno nello stesso gruppo, tendono a vivere perdite insieme.
Una possibile applicazione è selezionare un asset da ogni cluster: così si
evita di riempire il portafoglio con molti titoli che crollano nello stesso
modo. È un criterio di diversificazione, non una strategia completa né una
garanzia contro le crisi.

\section{Cosa dovrei portare all'incontro col relatore}

\subsection*{La spiegazione breve}
``Schmidt misura il limite della dipendenza estrema; Hua--Joe studia la velocità
quando il limite è zero. Durante et al. guardano la curva a soglie finite. La
TCF usa la diagonale della copula: sotto $0.5$ conta le perdite congiunte,
sopra $0.5$ i valori alti congiunti. Dividendo per la probabilità marginale
otteniamo una quota leggibile. La curva può restare alta a soglie osservabili
anche se il TDC finale è nullo. Gli autori usano questa curva per scegliere
copule con una mappa di distanze e per raggruppare mercati che tendono a perdere
insieme.''

\subsection*{Cosa devo aver capito}
\begin{enumerate}[leftmargin=*]
  \item La diagonale $\delta_C(t)=C(t,t)$ è una probabilità nella zona in cui
  due percentili superano o non superano la stessa soglia.
  \item Il TDC è un limite; la TCF è una curva a soglia finita.
  \item $t<0.5$ significa lower tail, $t>0.5$ upper tail; i denominatori $t$
  e $1-t$ trasformano un conteggio congiunto in una quota condizionata.
  \item Una TCF alta o lenta segnala più rischio congiunto osservabile.
  \item La stima empirica è un conteggio di punti nel quadrato dei percentili,
  quindi diventa rumorosa se il campione contiene pochi estremi.
  \item MDS è una mappa costruita da distanze tra curve; il piano non è una
  nuova variabile statistica.
  \item Nel clustering, i modelli marginali tolgono scala e volatilità; la
  copula misura la dipendenza tra percentili residui.
\end{enumerate}

\subsection*{Cosa non devo ancora saper fare}
Non devo ancora dimostrare i teoremi sulle diagonali, stimare da solo un modello
ARMA--GARCH completo o ricostruire tutta la teoria MDS. Devo però saper dire
perché servono: la diagonale rende visibile una parte della copula, ARMA--GARCH
prepara confronti equi tra serie, MDS visualizza distanze già calcolate.

\Salta per ora le condizioni tecniche per supporto pieno e copule assolutamente
continue, le costruzioni patchwork e i dettagli della funzione stress. Torna a
queste parti solo se la tesi diventa teorica sulle diagonali prescritte.

\subsection*{Domande concrete per il relatore}
\begin{description}[leftmargin=3.2cm,style=nextline]
  \item[Concettuali] Per la nostra applicazione è più importante il TDC, la
  TCF completa o il tail order? Quale coda rappresenta meglio le perdite?
  \item[Stima] Come scegliere $\varepsilon$ in funzione di $n$? Quanti punti
  estremi sono necessari per considerare leggibile una curva? Come costruire
  bande di incertezza per la TCF?
  \item[Copule] È sensato calibrare famiglie con la stessa tau centrale e poi
  confrontarle con TCF? Dovremmo pesare di più la lower tail nella distanza?
  \item[Serie finanziarie] ARMA--GARCH con innovazioni $t$ è adeguato per i
  nostri dati? Serve testare asimmetrie lower/upper prima del clustering?
  \item[Portafoglio] Il criterio ``un asset per cluster'' è solo esplorativo o
  può entrare in una procedura di selezione con vincoli di cardinalità?
  \item[Collegamento] Schmidt offre TDC, stima e VaR; Hua--Joe aggiunge ordine;
  Durante et al. aggiunge la TCF e strumenti grafici. Il relatore vede questi
  tre paper come una sequenza per la tesi o come alternative?
\end{description}

\subsection*{Legenda della notazione e pronuncia}
\begin{description}[leftmargin=3.8cm,style=nextline]
  \item[$C$] copula: come si combinano i percentili; si legge ``ci''.
  \item[$\delta_C(t)$] sezione diagonale $C(t,t)$; delta si legge ``delta''.
  \item[$\lambda_L,\lambda_U$] TDC lower e upper, cioè i limiti asintotici;
  lambda si legge ``lambda''.
  \item[$q_C(t)$] TCF: quota di estremi congiunti a una soglia finita; si può
  leggere ``cu di ti''.
  \item[$\beta_C$] associazione centrale di Blomqvist; beta si legge ``beta''.
  \item[$t$] percentile/soglia: sotto $0.5$ lower, sopra $0.5$ upper.
  \item[$F_n,G_n,C_n$] distribuzioni e copula empiriche ottenute dai ranghi.
  \item[$U_i,V_i$] percentili empirici o PIT delle osservazioni.
  \item[$d(C_i,C_j)$] distanza integrata tra due TCF; ``d'' significa distanza.
  \item[MDS] mappa bidimensionale costruita da una matrice di distanze.
  \item[$d_{ij}$] dissimilarità lower-tail tra asset $i$ e $j$.
  \item[$n,T$] rispettivamente numero di coppie in un campione e numero di
  giorni nella serie temporale.
  \item[$k$] indice della griglia $k/T$, lettera latina ``k'' (cappa), non il
  tail order $\kappa$ di Hua--Joe.
  \item[PIT] probability integral transform: trasformazione di un residuo nel
  suo percentile stimato.
  \item[ARMA--GARCH] modello per dinamica media e volatilità marginale, non per
  la copula tra asset.
\end{description}

\end{document}
